% ============================================================
%  Experiment 7 --- DC Motor Speed Measurement
% ============================================================
\experimentheader{7}{DC Motor Speed Measurement}{Instrumentation --- Measuring the Speed of a DC Motor Using a Computer}{Instrumentation Lab}

\begin{objectives}
\begin{enumerate}
  \item To understand the methods of acquiring measurements automatically with the aid of a computer.
  \item To become familiar with the data-acquisition (DAQ) environment.
  \item To become familiar with signal conditioning in a computer-based data-acquisition system.
  \item To design and build a signal-conditioning circuit whose parameters are justified by a measurement, rather than assumed.
  \item To implement a computer-based measurement system for the rotational speed of a DC motor.
\end{enumerate}
\end{objectives}

\begin{equipment}
\begin{itemize}
  \item DC motor trainer with a 12-slot wheel encoder
  \item Digital storage oscilloscope (DSO)
  \item Components for the student-designed filter and signal-conditioning circuit (op-amps, resistors, capacitors)
  \item ADAM-3968 wiring terminal board
  \item PCIE-1816 multifunction DAQ card, installed in the host PC
  \item PC with a C/C++ or Python development environment
\end{itemize}
\end{equipment}

\begin{introduction}
In most measurement processes it is often necessary to acquire data automatically, without a human operator. A computer-based instrument has several advantages over conventional measuring instruments: it can acquire measurements quickly and accurately, sustain real-time control, and operate for extended periods with consistent precision.

A computer-based measurement system typically contains the following elements (Figure~\ref{fig:exp7-daq-overview}): a transducer, a signal-conditioning stage, data-acquisition hardware, analysis hardware, and software running on the host PC.

\textbf{Signal conditioning} means manipulating an analogue signal so that it meets the requirements of the next processing stage --- most commonly, so that it is suitable for an analogue-to-digital converter. Signal conditioning can include amplification, filtering, range matching, and isolation.

\textbf{Filtering} is the most common signal-conditioning function: frequencies that do not carry valid data are attenuated. In this experiment the encoder produces a pulse train whose useful frequency content is bounded, and any filter used ahead of the DAQ card should be designed --- not assumed --- around that bound.

\textbf{Amplification} increases the resolution of the input signal and improves its signal-to-noise ratio.

\textbf{Isolation} allows a signal to pass from its source to the measurement device without a direct physical (galvanic) connection, protecting sensitive equipment from the signal source and vice versa.

\textbf{The Sallen--Key topology} is a widely used active-filter structure, valued for its simplicity: it realises a two-pole (12\,dB/octave) low-pass, high-pass, or band-pass response using a single op-amp configured as a near-unity-gain voltage-controlled voltage source (VCVS), together with two resistors and two capacitors that set the cutoff frequency and $Q$.
\end{introduction}

\begin{boxedfigure}
\centering
\begin{tikzpicture}[node distance=8mm and 10mm]
  \node[block] (trans) {Transducer};
  \node[block, right=of trans] (cond) {Signal\\Conditioning};
  \node[block, right=of cond] (daq) {Data\\Acquisition};
  \node[block, below=8mm of cond] (sw) {Software};
  \node[block, right=of sw] (pc) {Personal\\Computer};

  \draw[arr] (trans) -- (cond);
  \draw[arr] (cond) -- (daq);
  \draw[arr] (daq.south) |- (pc.north);
  \draw[arr] (sw) -- (pc);
\end{tikzpicture}
\captionof{figure}{Computer-based measurement system: transducer, signal conditioning, DAQ hardware, and analysis software.}
\label{fig:exp7-daq-overview}
\end{boxedfigure}

\begin{procedure}

\textbf{Activity 1 --- Characterising the encoder signal and designing the filter}\vspace{2pt}

\textit{Objective: measure the encoder's pulse characteristics and design a signal-conditioning filter around the measured cutoff frequency, rather than an assumed one.}

\begin{enumerate}
  \item Wire the DC motor trainer so that the 12-slot wheel encoder's output is available at a test point, and connect this test point to the oscilloscope.
  \item Run the motor across its usable speed range and observe the encoder pulse train on the oscilloscope. Each revolution produces 12 pulses (one per slot).
  \item At a representative operating speed, measure the pulse period and estimate the corresponding fundamental frequency and its dominant harmonics directly from the oscilloscope trace.
  \item From this measurement, determine a suitable cutoff frequency for an anti-aliasing / noise-rejection low-pass filter: the cutoff must pass the pulse's fundamental frequency (across the motor's full speed range) while rejecting high-frequency switching noise.
  \item \textbf{Design task:} using the Sallen--Key topology introduced above, design a second-order active low-pass filter with the cutoff frequency determined in step 4. Select a topology (unity-gain or gain-set), calculate resistor and capacitor values, and justify your component choices in your report.
  \item Build and test the filter on its own (e.g.\ with a function generator) to verify its measured cutoff frequency and roll-off match your design before connecting it to the encoder signal.
  \item \textbf{Design task:} design the remaining signal-conditioning stage needed to present the filtered encoder pulses to the PCIE-1816's digital input with valid logic levels (buffering, level-shifting, or pulse-shaping as required). Justify your design choices.
\end{enumerate}

\begin{needsupdate}
The filter and signal-conditioning circuit are student design deliverables: there is no fixed cutoff frequency or component value specified here. Document your oscilloscope measurements, design calculations, and any deviation between the designed and measured filter response in your report.
\end{needsupdate}

\textbf{Activity 2 --- Interfacing to the DAQ hardware}\vspace{2pt}

\textit{Objective: connect the conditioned encoder signal to the PC via the ADAM-3968 and PCIE-1816.}

\begin{enumerate}
  \item Connect the output of your signal-conditioning circuit to a digital input (or counter/timer input, if using the PCIE-1816's onboard counter) on the ADAM-3968 terminal board.
  \item Connect the ADAM-3968 to the PCIE-1816 DAQ card installed in the host PC via the supplied cable.
  \item Verify the wiring using the DAQ card's bundled configuration/monitoring utility before writing any code, to confirm the expected pulses are being registered.
\end{enumerate}

\begin{boxedfigure}
\centering
\resizebox{0.98\linewidth}{!}{%
\begin{tikzpicture}[node distance=12mm and 16mm]
  \node[block, minimum width=2.6cm] (trainer) {DC Motor\\Trainer};
  \node[block, right=of trainer, minimum width=2.8cm] (filt) {Filter \&\\Signal Cond.\\\scriptsize(student design)};
  \node[block, right=of filt, minimum width=2.4cm] (adam) {ADAM-3968};
  \node[block, right=of adam, minimum width=2.6cm] (pcie) {PCIE-1816};
  \node[block, right=of pcie, minimum width=2.4cm] (pc) {PC\\\scriptsize C/C++/Python GUI};

  \draw[arr] (trainer) -- (filt) node[midway, above, font=\tiny] {encoder pulses};
  \draw[arr] (filt) -- (adam);
  \draw[arr] (adam) -- (pcie);
  \draw[arr] (pcie) -- (pc);
\end{tikzpicture}
}
\captionof{figure}{Signal path from the DC motor trainer to the host PC.}
\label{fig:exp7-signalpath}
\end{boxedfigure}

\textbf{Activity 3 --- Speed measurement software}\vspace{2pt}

\textit{Objective: implement software to compute and display motor speed in real time.}

The rotational speed can be found from the pulse period, using the number of pulses produced per revolution (12, for this encoder):
\[
\text{Period} = \frac{t_{\text{stop}} - t_{\text{start}}}{\text{CLK\_TCK}}, \qquad
\omega = \frac{2\pi}{12 \times \text{Period}} \ \left[\text{rad/s}\right]
\]
where $t_{\text{start}}$ and $t_{\text{stop}}$ mark successive encoder edges and $\text{CLK\_TCK}$ is the counter's clock tick rate.

\begin{boxedfigure}
\centering
\begin{tikzpicture}[node distance=7mm and 10mm]
  \node[startstop] (start) {Start};
  \node[block, below=of start] (init) {Initialise DAQ card};
  \node[block, below=of init] (read) {Acquire encoder pulse timing};
  \node[block, below=of read] (calc) {Compute period and $\omega$ (rad/s)};
  \node[block, below=of calc] (disp) {Display / log result};
  \node[decision, below=of disp, minimum width=2.6cm] (dec) {End of\\measurement?};
  \node[startstop, below=of dec] (end) {End};

  \draw[arr] (start) -- (init);
  \draw[arr] (init) -- (read);
  \draw[arr] (read) -- (calc);
  \draw[arr] (calc) -- (disp);
  \draw[arr] (disp) -- (dec);
  \draw[arr] (dec) -- (end) node[midway, right, font=\small] {Yes};
  \draw[arr] (dec.east) -- ++(2.4,0) |- (read.east) node[pos=0.25, above, font=\small] {No};
\end{tikzpicture}
\captionof{figure}{Measurement loop for the speed-acquisition software.}
\label{fig:exp7-flowchart}
\end{boxedfigure}

\begin{enumerate}
  \item In C, C++, or Python, write a program that:
  \begin{itemize}
    \item initialises the PCIE-1816 (via the vendor driver / SDK, or the DAQ card's counter API),
    \item times successive encoder pulses (or reads the onboard counter) to obtain the pulse period,
    \item computes the rotational speed in rad/s (or rpm) using the relation above, and
    \item displays the running value in a simple GUI, updating in real time.
  \end{itemize}
  \item Interface your filter/conditioning circuit, the ADAM-3968, and the PCIE-1816 together and perform measurements across a range of motor input voltages. Compare your computed speed against a direct oscilloscope measurement of the encoder period at the same operating point.
  \item Tabulate your results (Table~\ref{tab:exp7-results}).
  \item Plot computed speed (from your software) against the oscilloscope-measured speed.
  \item Discuss and comment on your results, including any discrepancy between the two measurement methods and possible causes (filter phase delay, counting resolution, etc.).
\end{enumerate}

\begin{boxedtable}
\centering
\begin{tabular}{lccc}
\toprule
& \multicolumn{3}{c}{Input voltage} \\
\cmidrule(lr){2-4}
& \# 1 & \# 2 & \# 3 \\
\midrule
Speed (rad/s) --- computer & & & \\
Speed (rad/s) --- oscilloscope & & & \\
\bottomrule
\end{tabular}
\captionof{table}{Computed vs.\ oscilloscope-measured motor speed.}
\label{tab:exp7-results}
\end{boxedtable}

\end{procedure}

\begin{reportq}
\begin{enumerate}
  \item Present your oscilloscope measurements of the encoder pulse and your resulting filter design (topology, calculations, and component values), together with the measured response of the built filter.
  \item Present your design of the signal-conditioning stage between the filter and the PCIE-1816 digital input, with justification.
  \item Compare your software-computed speed with the oscilloscope-measured speed across your tested operating points, and discuss the sources of any discrepancy.
  \item What is the theoretical maximum measurable speed of this system, given the encoder resolution (12 slots/revolution) and your software's timing resolution?
\end{enumerate}
\end{reportq}
